\section{The transfer is an all-pass filter}\label{sec:allpass}

Section~\ref{sec:generator} wrote the generator out. This section identifies the
transfer, and the identification is the hinge of the paper: one elementary
proposition, and four consequences that occupy the rest of it.

\subsection{The symbol is unimodular}

\begin{proposition}[Unimodularity]\label{prop:unimodular}
For every real $\tau$ and every $\omega$ with $\xi(\frac12+\omega+i\tau)\neq0$,
\[
 |K_\omega(i\tau)|=1 .
\]
No hypothesis is used beyond the functional equation and the reality of $\xi$ on
the real axis.
\end{proposition}

\begin{proof}
$\xi(s)=\xi(1-s)$ gives
$\xi(\frac12-\omega+i\tau)=\xi(\frac12+\omega-i\tau)$, and $\xi$ has real Taylor
coefficients about $s=\frac12$, so $\xi(\bar s)=\overline{\xi(s)}$ and
$\xi(\frac12+\omega-i\tau)=\overline{\xi(\frac12+\omega+i\tau)}$. Numerator and
denominator of \eqref{eq:symbol} are complex conjugates.
\end{proof}

So $\Theta_\omega$ is an \emph{all-pass filter}: it changes no amplitude at any
frequency, only phase. Writing
\begin{equation}\label{eq:phase}
 \xi(\tfrac12+\omega+i\tau)=R(\tau,\omega)e^{i\Psi(\tau,\omega)},
 \qquad K_\omega(i\tau)=e^{-2i\Psi(\tau,\omega)},
\end{equation}
the entire content of the transfer on the critical line is the single real
function $\Psi$. Everything below is a statement about where that phase goes.

\begin{remark}[A trap: which contour may be moved]\label{rem:trap}
Unimodularity forces
$\re a_\omega(i\tau)=-\partial_\omega\log|K_\omega(i\tau)|=0$, which appears to
make $Q_{\omega,L}$ vanish. It does not, and the reason is worth recording
because the two symbols behave differently.

$K_\omega$ has poles only where $\xi(\frac12+p+\omega)=0$, at
$p=\rho-\frac12-\omega$. Under the Riemann hypothesis those sit on
$\re p=-\omega$, so the inverse Laplace contour \emph{may} be moved from
$\re p=\eta>1$ to $\re p=0$. But $a_\omega$ has poles at
$p=\rho-\frac12\pm\omega$, hence also on $\re p=+\omega$, inside the right
half-plane, so for it the contour \emph{may not} be moved. The continuation
$a_\omega(i\tau)$ is therefore not the symbol of the causal operator, and the
residues crossed in the shift are the whole of the form: as $\omega\downarrow0$
they become $2\pi\sum_\rho\delta(\tau-\gamma_\rho)$, which is
\eqref{eq:spectral-target}.

Relatedly, the four apparent poles at $p=\pm\frac12\pm\omega$ in the factor
decomposition \eqref{eq:threefactors} \emph{cancel} between the three factors:
$\psi$ has a pole at $p=-\frac12$ and $\zeta'/\zeta$ one at $p=\frac12$, with
matching residues. The proof of Theorem~\ref{thm:kernel} is unaffected, because
it transforms each factor on $\re p=\eta>1$ where all converge and sums; but
$a_\omega$ does not have those poles.
\end{remark}

\subsection{Consequence one: the Weil symbol is a group delay}

\begin{proposition}[Group delay]\label{prop:groupdelay}
The group delay of the transfer is
\[
 -\frac{d}{d\tau}\arg K_\omega(i\tau)=2\,\partial_\tau\Psi
 =2\,\re\frac{\xi'}{\xi}\Big(\tfrac12+\omega+i\tau\Big),
\]
and $\partial_\tau\Psi$ is $\pi$ times the density of zeros. Its smooth part is
$\theta'(\tau)$, $\theta$ the Riemann--Siegel theta function, so the group delay
is $2\theta'(\tau)=b(\tau^2)+w_0$.
\end{proposition}

\begin{proof}
The first equality is \eqref{eq:phase}; the second is
$\partial_\tau\arg\xi=\re(\xi'/\xi)$ along the vertical line. That
$\partial_\tau\arg\xi(\frac12+\omega+i\tau)$ is $\pi$ times the zero density is
the argument principle, and $N(T)=\theta(T)/\pi+1+S(T)$ identifies the smooth
part. The final equality is the density identity of \cite{Selection}.
\end{proof}

The companion investigation obtained $b(\tau^2)+w_0=2\theta'(\tau)$ as an identity
between special functions and drew from it that the contact constant "is not an
independent object requiring a mechanism". Proposition~\ref{prop:groupdelay} says
why: the archimedean symbol is a \emph{phase derivative}, and a phase has no
separate ingredients. It also reads the unboundedness rule of \cite{Selection} as
a statement about time rather than about divergence: the delay grows like
$\log\tau$, so a packet at frequency $N$ is held for $\sim\log N$, which is
$Q_L[\chi e^{iNx}]/\log N\to\lVert\chi\rVert^2$.

\subsection{Consequence two: the zeros are resonances}

\begin{proposition}[Resonant line shape]\label{prop:resonance}
As $\omega\downarrow0$, near each $\gamma_\rho$,
\[
 |K_\omega(i\tau)-1|^2=\frac{4\omega^2}{\omega^2+(\tau-\gamma_\rho)^2}+o(1),
\]
a Lorentzian of height $4$ and half-width $\omega$, of mass $4\pi\omega$.
Consequently
\begin{equation}\label{eq:absorption}
 \frac1{2\pi}\int_{\R}|K_\omega(i\tau)-1|^2|\widehat F(\tau)|^2\dd\tau
 \;\longrightarrow\;2\omega\sum_\rho|\widehat F(\gamma_\rho)|^2
 \;=\;2\omega\,Q_{0,L}[f].
\end{equation}
\end{proposition}

\begin{proof}
Near a simple zero, $\xi(\frac12+\omega+i\tau)\approx\xi'(\rho)(\omega+i(\tau-\gamma))$,
and by Proposition~\ref{prop:unimodular} the numerator is its conjugate, so
$K_\omega(i\tau)$ is a Blaschke factor times a constant phase. Requiring
$K_\omega\to1$ away from the zero as $\omega\downarrow0$ fixes that phase, leaving
$K_\omega(i\tau)\approx\big(\omega-i(\tau-\gamma)\big)/\big(\omega+i(\tau-\gamma)\big)$,
for which $|K_\omega-1|^2$ is the stated Lorentzian exactly; its mass is
$\int4\omega^2(\omega^2+u^2)^{-1}\dd u=4\pi\omega$. Summing over the zeros and
dividing by $2\pi$ gives \eqref{eq:absorption}, whose right side is
\eqref{eq:spectral-target}.
\end{proof}

So the leading behaviour of the defect,
$\langle f,D_{\omega,L}f\rangle=2\omega Q_{0,L}[f]+o(\omega)$, is not a formal
consequence of the flow identities but a statement about line shapes: each zero
is a resonance of width $\omega$ and fixed integrated strength, and the Weil form
is the total resonant response of the input. Appendix~\ref{app:checks} verifies
the line shape, its mass and the window sum in a model built from published zero
ordinates, where all three hold exactly.

\subsection{Consequence three: under RH the defect is a boundary flux}

\begin{proposition}[Inner, unitary, and the flux identity]\label{prop:flux}
Assume the Riemann hypothesis. Then $K_\omega$ is an inner function of the right
half-plane, with zeros exactly at $p=\omega+i\gamma_\rho$; $V_\omega$ is unitary
on $L^2(\R)$ and causal; and for $f$ supported in $I_L$,
\begin{equation}\label{eq:flux}
 \langle f,D_{\omega,L}f\rangle=\lVert(1-P_L)V_\omega f\rVert^2
 =\int_{L/2}^{\infty}\big|(V_\omega f)(x)\big|^2\dd x .
\end{equation}
\end{proposition}

\begin{proof}
Under RH the poles of $K_\omega$ lie on $\re p=-\omega$, so $K_\omega$ is analytic
on $\re p\geq0$; it is bounded there, the $\Gamma$-ratio giving
$K_\omega(p)=O(|p|^{-\omega})$ while the $\zeta$ and pole ratios tend to $1$; and
it is unimodular on the boundary by Proposition~\ref{prop:unimodular}. Hence it is
inner, the Blaschke condition $\sum_\rho\omega/(1+\gamma_\rho^2)<\infty$ holding.
No pole lies between $\re p=\eta$ and $\re p=0$, so the inverse Laplace contour
may be moved to the imaginary axis and $V_\omega$ is the Fourier multiplier by
$K_\omega(i\cdot)$, of modulus one, hence unitary. Then
$\lVert f\rVert^2-\lVert P_LV_\omega f\rVert^2
=\lVert V_\omega f\rVert^2-\lVert P_LV_\omega f\rVert^2
=\lVert(1-P_L)V_\omega f\rVert^2$, and causality gives
$\supp V_\omega f\subseteq[-L/2,\infty)$, so all the escaping mass is to the
right.
\end{proof}

\eqref{eq:flux} is the sharpest statement in the paper. The contraction defect is
not an abstract operator quantity: it is the \textbf{flux of $|V_\omega f|^2$
through the leading endpoint of the interval}, and Weil positivity on $I_L$ says
the flow pushes mass out of the moving end and never pulls it in. Section~\ref{sec:endpoint}
finds the same point from the geometry, as the location of a rank-one variation.

\begin{remark}[Direction of the hypothesis]\label{rem:direction}
Proposition~\ref{prop:flux} uses RH; it is the easy direction of the program made
explicit, not a route to a proof. Its value is that it exhibits the mechanism,
which Sections~\ref{sec:region} and \ref{sec:endpoint} then use.
\end{remark}

\subsection{The scattering dictionary}

Collected, Propositions~\ref{prop:unimodular}--\ref{prop:flux} are a definition.

\begin{center}
\begin{tabular}{@{}p{0.46\textwidth}p{0.46\textwidth}@{}}
\toprule
Established here & Scattering name\\
\midrule
$k_\omega$ supported in $[0,\infty)$; $K_\omega$ analytic in a right half-plane &
causality, via Paley--Wiener\\
$|K_\omega(i\tau)|=1$, exactly and unconditionally &
unitarity on the energy shell\\
poles at $-\omega+i\gamma_\rho$ (Remark~\ref{rem:trap}) &
resonances of width $\omega$\\
$-\,d\varphi/d\tau=2\theta'(\tau)$ (Prop.~\ref{prop:groupdelay}) &
Wigner--Smith time delay $=$ density of states \cite{Smith}\\
$\Psi/\pi\to N(\tau)$ (Prop.~\ref{prop:groupdelay}) &
Krein spectral shift function \cite{BirmanKrein}\\
\bottomrule
\end{tabular}
\end{center}

Two further points make the identification specific. $K_\omega$ is a
\emph{ratio} --- $\xi$ at the reflected argument over $\xi$ at the argument, with
the functional equation supplying the reflection --- and a unimodular ratio of that
shape is what a one-channel \emph{reflection coefficient} is; a half-infinite line
with a boundary condition at one end has exactly one scattering channel and its
entire $S$-matrix is such a coefficient. And the argument-principle bookkeeping of
Proposition~\ref{prop:groupdelay} is Levinson's theorem, the total phase change
per resonance being $-2\pi$.

\begin{remark}[Scope of this reading]\label{rem:scatteringscope}
Reading $\tau$ as an energy and $x$ as a scattering coordinate is an
identification the program does not make; the caveat attached to the channel-tower
rule of \cite{Selection} applies here in full, and Section~\ref{sec:conditions}
gives the first motivated reason to make it without removing it. The
neighbourhood is also crowded: zeros-as-resonances is old. Nothing here claims
the idea. What is this program's own is that the object arrives from its own
shifted transfer, with the constants fixed --- the delay is $2\theta'$ exactly and
the resonance widths are exactly $\omega$.
\end{remark}
